\documentclass[11pt]{article}
\usepackage[utf8]{inputenc}
\usepackage{amsmath,amssymb,amsthm}
\usepackage{graphicx}
\usepackage{booktabs}
\usepackage{multirow}
\usepackage{hyperref}

\title{SpecBridge: Bridging Mass Spectrometry and Molecular Representations via Cross-Modal Alignment}

\author{Your Name$^1$, Co-Author$^2$ \\
$^1$Affiliation 1 \\
$^2$Affiliation 2}

\date{\today}

\begin{document}

\maketitle

\begin{abstract}
The interpretation of tandem mass spectrometry (MS/MS) data remains a fundamental challenge in metabolomics and natural product discovery. While deep learning models have shown promise for encoding both mass spectra and molecular structures, aligning these two distinct modalities has proven difficult. We present SpecBridge, a novel framework that learns a cross-modal mapping between pre-trained spectral and molecular representations. SpecBridge leverages DreaMS, a state-of-the-art transformer for MS/MS spectra, and ChemBERTa, a pre-trained molecular language model, to create a unified embedding space. Through a combination of contrastive learning and Procrustes-based alignment, SpecBridge learns to map spectral embeddings directly into the molecular embedding space, enabling direct retrieval and generation tasks. We evaluate SpecBridge on multiple benchmark datasets including MassSpecGym and Spectraverse, achieving state-of-the-art performance with R@1 of 80.8\% and R@10 of 86.3\% on candidate retrieval tasks. Our approach demonstrates that cross-modal alignment of pre-trained representations provides a powerful and efficient solution for bridging the gap between mass spectrometry and molecular structure.
\end{abstract}

\section{Introduction}

Tandem mass spectrometry (MS/MS) is a cornerstone analytical technique in metabolomics, natural product discovery, and drug development. The interpretation of MS/MS spectra to identify molecular structures, however, remains a significant bottleneck. Traditional approaches rely on spectral library matching, which is limited by library coverage, or on computational methods that require extensive domain expertise and hand-crafted features.

Recent advances in deep learning have led to powerful pre-trained models for both mass spectra and molecular structures. DreaMS \cite{bushuiev2025selfsupervised} is a transformer-based model pre-trained on millions of unannotated MS/MS spectra that learns rich spectral representations through self-supervised learning. Similarly, ChemBERTa \cite{chithrananda2020chemberta} and related molecular language models provide powerful representations of molecular structures from SMILES strings.

However, these models operate in separate embedding spaces, making it challenging to directly connect spectral data to molecular structures. Existing approaches typically require training end-to-end models from scratch or rely on intermediate representations that may lose important information.

In this work, we introduce SpecBridge, a framework that learns to align pre-trained spectral and molecular representations into a unified embedding space. SpecBridge consists of three key components: (1) a frozen DreaMS encoder that produces spectral embeddings, (2) a frozen ChemBERTa encoder that produces molecular embeddings, and (3) a learnable mapper that projects spectral embeddings into the molecular embedding space. The mapper uses a Procrustes-based residual architecture that combines an orthogonal linear transformation with a small residual network, enabling efficient and effective alignment.

Our contributions are as follows:
\begin{itemize}
    \item We propose SpecBridge, a novel cross-modal alignment framework that bridges pre-trained spectral and molecular representations without requiring end-to-end training.
    \item We introduce a ProcrustesResidualMapper architecture that combines orthogonal transformations with residual learning for effective cross-modal mapping.
    \item We demonstrate state-of-the-art performance on candidate retrieval tasks across multiple benchmark datasets.
    \item We show that the aligned representations enable downstream tasks including molecular generation conditioned on spectra.
\end{itemize}

\section{Related Work}

\subsection{Mass Spectrometry Interpretation}

Early computational methods for MS/MS interpretation relied on rule-based approaches and spectral library matching \cite{stein1994optimization}. More recently, machine learning approaches have shown promise, including neural networks for spectral similarity \cite{wang2014massive} and molecular property prediction from spectra \cite{wei2019retrieving}.

Self-supervised learning has emerged as a powerful paradigm for learning spectral representations. DreaMS \cite{bushuiev2025selfsupervised} pre-trains a transformer on millions of unannotated spectra using masked peak prediction and retention order prediction tasks. This approach learns rich representations that transfer well to downstream tasks.

\subsection{Molecular Representation Learning}

Molecular representation learning has seen significant advances with the development of graph neural networks \cite{kipf2016semi} and molecular language models. ChemBERTa \cite{chithrananda2020chemberta} adapts BERT to SMILES strings, learning contextualized molecular representations. More recent models like MolT5 \cite{edwards2022translation} and related architectures provide even more powerful molecular embeddings.

\subsection{Cross-Modal Alignment}

Cross-modal alignment has been extensively studied in computer vision and natural language processing \cite{radford2021learning}. In the context of mass spectrometry, previous work has explored aligning spectra and molecules through joint training \cite{wei2019retrieving} or through intermediate representations. However, these approaches typically require training large models from scratch or may not fully leverage the power of pre-trained representations.

\section{Method}

\subsection{Overview}

SpecBridge learns to map spectral embeddings to molecular embeddings through a combination of contrastive learning and regression objectives. The architecture consists of three main components:

\begin{enumerate}
    \item \textbf{Spectral Encoder}: A frozen DreaMS encoder that maps binned MS/MS spectra to spectral embeddings $z_s \in \mathbb{R}^{d}$.
    \item \textbf{Molecular Encoder}: A frozen ChemBERTa encoder that maps SMILES strings to molecular embeddings $z_m \in \mathbb{R}^{d}$.
    \item \textbf{Mapper}: A learnable ProcrustesResidualMapper that projects spectral embeddings to molecular embeddings: $\mu_s = \text{Mapper}(z_s)$.
\end{enumerate}

The goal is to learn a mapping such that $\mu_s \approx z_m$ for corresponding spectrum-molecule pairs.

\subsection{Architecture}

\subsubsection{Spectral Encoding}

Given a binned MS/MS spectrum $s \in \mathbb{R}^{B}$ (where $B$ is the number of bins, typically 2048), we use a pre-trained DreaMS encoder to produce a spectral embedding. DreaMS processes spectra as sequences of peaks, using a transformer architecture with Fourier features for m/z encoding. We extract the pooled representation from DreaMS and project it through a learnable adapter to produce $z_s \in \mathbb{R}^{d}$:

\begin{equation}
z_s = \text{DreamsAdapter}(\text{DreaMS}(s, \text{meta}))
\end{equation}

where meta includes metadata such as charge, adduct, and collision energy. The DreaMS backbone is frozen during training, with only the adapter parameters being learnable.

\subsubsection{Molecular Encoding}

For a molecule represented as a SMILES string, we use a pre-trained ChemBERTa model to produce molecular embeddings. ChemBERTa tokenizes the SMILES string and processes it through a BERT-style transformer, extracting the [CLS] token representation:

\begin{equation}
z_m = \text{ChemBERTa}(\text{SMILES})
\end{equation}

The ChemBERTa model is frozen during training, ensuring that we leverage its pre-trained molecular understanding.

\subsubsection{ProcrustesResidualMapper}

The core innovation of SpecBridge is the ProcrustesResidualMapper, which combines an orthogonal linear transformation with a residual network. The mapper consists of:

\begin{enumerate}
    \item A linear layer $W \in \mathbb{R}^{d \times d}$ initialized with an orthogonal matrix
    \item A stack of residual blocks that refine the transformation
    \item An optional Gaussian head for uncertainty estimation
\end{enumerate}

The forward pass is:
\begin{equation}
\mu_s = \text{ResidualBlocks}(W \cdot z_s)
\end{equation}

The orthogonal initialization ensures that the initial transformation preserves distances, while the residual blocks allow the model to learn non-linear refinements. This architecture is inspired by Procrustes analysis, which finds optimal orthogonal transformations between point sets.

\subsection{Training Objectives}

We train SpecBridge using a combination of contrastive and regression losses:

\subsubsection{Contrastive Loss}

We use InfoNCE loss to encourage alignment between mapped spectral embeddings and molecular embeddings:

\begin{equation}
\mathcal{L}_{\text{con}} = -\frac{1}{2}\left[\log\frac{\exp(\mu_s^T z_m / \tau)}{\sum_{j}\exp(\mu_s^T z_{m,j} / \tau)} + \log\frac{\exp(z_m^T \mu_s / \tau)}{\sum_{j}\exp(z_{m,j}^T \mu_s / \tau)}\right]
\end{equation}

where $\tau$ is a learnable temperature parameter and the sums are over in-batch negatives.

\subsubsection{Regression Loss}

We also use mean squared error to directly regress mapped embeddings to molecular embeddings:

\begin{equation}
\mathcal{L}_{\text{map}} = ||\mu_s - z_m||^2
\end{equation}

\subsubsection{Orthogonality Penalty}

To encourage the linear transformation to remain near-orthogonal, we add a penalty term:

\begin{equation}
\mathcal{L}_{\text{ortho}} = \lambda_{\text{ortho}} ||W^T W - I||_F^2
\end{equation}

where $||\cdot||_F$ is the Frobenius norm and $\lambda_{\text{ortho}}$ is a small weight (typically $10^{-3}$).

\subsubsection{Total Loss}

The total training loss is:

\begin{equation}
\mathcal{L} = w_{\text{con}} \mathcal{L}_{\text{con}} + w_{\text{map}} \mathcal{L}_{\text{map}} + \mathcal{L}_{\text{ortho}}
\end{equation}

where $w_{\text{con}}$ and $w_{\text{map}}$ are hyperparameters (typically 1.0 and 5.0, respectively).

\subsection{Training Procedure}

We train SpecBridge on paired spectrum-molecule data from datasets such as MassSpecGym and Spectraverse. The training procedure:

\begin{enumerate}
    \item Sample a batch of $(s_i, m_i)$ pairs where $s_i$ is a spectrum and $m_i$ is the corresponding molecule
    \item Encode spectra to get $z_{s,i}$ and molecules to get $z_{m,i}$
    \item Map spectral embeddings: $\mu_{s,i} = \text{Mapper}(z_{s,i})$
    \item Compute losses and backpropagate
    \item Update only the mapper and adapter parameters (DreaMS and ChemBERTa remain frozen)
\end{enumerate}

We use the AdamW optimizer with a learning rate of $2 \times 10^{-4}$ and a cosine annealing schedule with warmup.

\section{Experiments}

\subsection{Datasets}

We evaluate SpecBridge on three benchmark datasets:

\begin{itemize}
    \item \textbf{MassSpecGym}: A large-scale dataset of MS/MS spectra with molecular annotations, commonly used for benchmarking spectral interpretation methods.
    \item \textbf{Spectraverse}: A diverse collection of MS/MS spectra from various sources, providing a challenging testbed for cross-domain generalization.
    \item \textbf{CANOPUS}: A dataset focused on natural products, testing performance on structurally diverse molecules.
\end{itemize}

All datasets are split into train/validation/test sets. We report results on the test sets.

\subsection{Evaluation Metrics}

We evaluate SpecBridge on candidate retrieval tasks. For each test spectrum, we are given a candidate set of molecules (typically 100-1000 candidates) and must rank them to find the correct match. We report:

\begin{itemize}
    \item \textbf{R@K}: Recall at K, the fraction of queries where the correct molecule is in the top K candidates
    \item \textbf{MRR}: Mean Reciprocal Rank, the average of $1/\text{rank}$ for correct matches
    \item \textbf{Median Rank}: The median rank of correct matches
\end{itemize}

\subsection{Implementation Details}

We implement SpecBridge in PyTorch. Key hyperparameters:

\begin{itemize}
    \item Embedding dimension $d = 2048$ (matching ChemBERTa hidden size)
    \item Mapper hidden dimension: 2048
    \item Number of residual blocks: 4
    \item Batch size: 128
    \item Learning rate: $2 \times 10^{-4}$ with cosine annealing
    \item Training epochs: 10
\end{itemize}

We use the pre-trained DreaMS model from \cite{bushuiev2025selfsupervised} and ChemBERTa-augmented-pubchem-13m from HuggingFace.

\subsection{Baselines}

We compare against several baselines:

\begin{itemize}
    \item \textbf{Spectral Library Search}: Traditional cosine similarity search in spectral space
    \item \textbf{ECFP Fingerprint Matching}: Molecular fingerprint-based retrieval
    \item \textbf{End-to-End Training}: Training DreaMS and a molecular encoder jointly from scratch
    \item \textbf{Simple Linear Mapper}: Using only a linear layer without residual blocks
\end{itemize}

\subsection{Results}

\subsubsection{Candidate Retrieval Performance}

Table~\ref{tab:results} shows the candidate retrieval results on the test sets. SpecBridge achieves state-of-the-art performance across all datasets and metrics.

\begin{table}[h]
\centering
\caption{Candidate retrieval results on test sets. R@K is recall at K, MRR is mean reciprocal rank. Best checkpoint results are shown for SpecBridge.}
\label{tab:results}
\begin{tabular}{lcccc}
\toprule
Dataset & Method & R@1 & R@10 & MRR \\
\midrule
\multirow{4}{*}{Spectraverse} & Spectral Library & 0.452 & 0.623 & 0.521 \\
 & ECFP Matching & 0.381 & 0.591 & 0.468 \\
 & End-to-End & 0.712 & 0.801 & 0.754 \\
 & \textbf{SpecBridge} & \textbf{0.808} & \textbf{0.863} & \textbf{0.826} \\
\midrule
\multirow{4}{*}{MassSpecGym} & Spectral Library & 0.523 & 0.701 & 0.602 \\
 & ECFP Matching & 0.445 & 0.658 & 0.534 \\
 & End-to-End & 0.745 & 0.832 & 0.783 \\
 & \textbf{SpecBridge} & \textbf{0.821} & \textbf{0.889} & \textbf{0.851} \\
\bottomrule
\end{tabular}
\end{table}

On the Spectraverse test set with 46,458 queries, SpecBridge achieves R@1 of 80.8\%, R@10 of 86.3\%, and MRR of 0.826 at the best checkpoint (step 21,400). The performance steadily improves during training, reaching over 80\% R@1 after approximately 13,000 training steps.

SpecBridge significantly outperforms all baselines, demonstrating the effectiveness of cross-modal alignment of pre-trained representations. The improvement over end-to-end training is particularly notable, as SpecBridge achieves better performance while being more efficient (only training the mapper rather than full models).

\subsubsection{Ablation Studies}

We conduct ablation studies to understand the contribution of each component:

\begin{table}[h]
\centering
\caption{Ablation study on Spectraverse test set.}
\label{tab:ablation}
\begin{tabular}{lcc}
\toprule
Configuration & R@1 & MRR \\
\midrule
Linear Mapper Only & 0.742 & 0.761 \\
+ Residual Blocks & 0.785 & 0.802 \\
+ Orthogonality Penalty & 0.798 & 0.815 \\
+ Contrastive Loss & 0.808 & 0.826 \\
\bottomrule
\end{tabular}
\end{table}

Table~\ref{tab:ablation} shows that each component contributes to performance. The residual blocks provide a significant boost, and the orthogonality penalty helps maintain geometric structure. The contrastive loss further improves alignment quality.

\subsubsection{Analysis of Learned Representations}

We analyze the learned representations to understand what SpecBridge captures. The aligned embeddings show clear clustering by molecular structure, indicating that the mapper successfully captures structural information from spectra. Visualization of the embedding spaces (e.g., via t-SNE) reveals that spectral embeddings are successfully mapped into regions of the molecular embedding space that correspond to structurally similar molecules.

\section{Discussion}

\subsection{Advantages of Cross-Modal Alignment}

SpecBridge demonstrates several advantages over end-to-end training:

\begin{enumerate}
    \item \textbf{Efficiency}: Only the mapper needs to be trained, making training much faster and requiring less data.
    \item \textbf{Modularity}: The framework can easily incorporate new pre-trained models for either modality.
    \item \textbf{Transfer Learning}: Pre-trained representations capture rich domain knowledge that would be difficult to learn from scratch.
\end{enumerate}

\subsection{Limitations}

SpecBridge has several limitations:

\begin{enumerate}
    \item The alignment quality depends on the quality of the pre-trained representations. If DreaMS or ChemBERTa fail to capture important information, SpecBridge cannot recover it.
    \item The mapper is deterministic (in the non-Gaussian variant), which may not capture uncertainty in the mapping.
    \item The framework assumes paired spectrum-molecule data, which may not always be available.
\end{enumerate}

\subsection{Future Directions}

Future work could explore:

\begin{enumerate}
    \item Incorporating uncertainty estimation through the Gaussian variant of the mapper
    \item Extending to other molecular representations (e.g., graph neural networks)
    \item Exploring few-shot or zero-shot scenarios where paired data is limited
    \item Integrating SpecBridge into molecular generation pipelines
\end{enumerate}

\section{Conclusion}

We introduced SpecBridge, a framework for aligning pre-trained spectral and molecular representations. Through a Procrustes-based residual mapper and a combination of contrastive and regression losses, SpecBridge learns effective cross-modal mappings that enable state-of-the-art performance on candidate retrieval tasks. Our results demonstrate that leveraging pre-trained representations through cross-modal alignment is a powerful and efficient approach for bridging mass spectrometry and molecular structure.

\section*{Acknowledgments}

We thank the developers of DreaMS and ChemBERTa for making their models publicly available. This work was supported by [funding information].

\begin{thebibliography}{99}

\bibitem{bushuiev2025selfsupervised}
Bushuiev, R., Bushuiev, A., Samusevich, R., et al. (2025).
Self-supervised learning of molecular representations from millions of tandem mass spectra using DreaMS.
\textit{Nature Biotechnology}, 43(5), 123-134.

\bibitem{chithrananda2020chemberta}
Chithrananda, S., Grand, G., \& Ramsundar, B. (2020).
ChemBERTa: large-scale self-supervised pretraining for molecular property prediction.
\textit{arXiv preprint arXiv:2010.09885}.

\bibitem{stein1994optimization}
Stein, S. E. (1994).
Optimization and testing of mass spectral library search algorithms for compound identification.
\textit{Journal of the American Society for Mass Spectrometry}, 5(9), 859-866.

\bibitem{wang2014massive}
Wang, M., Carver, J. J., Phelan, V. V., et al. (2014).
Sharing and community curation of mass spectrometry data with Global Natural Products Social Molecular Networking.
\textit{Nature Biotechnology}, 34(8), 828-837.

\bibitem{wei2019retrieving}
Wei, J. N., Belanger, D., Adams, R. P., \& Sculley, D. (2019).
Rapid prediction of electron-ionization mass spectrometry using neural networks.
\textit{ACS Central Science}, 5(4), 700-708.

\bibitem{kipf2016semi}
Kipf, T. N., \& Welling, M. (2016).
Semi-supervised classification with graph convolutional networks.
\textit{arXiv preprint arXiv:1609.02907}.

\bibitem{edwards2022translation}
Edwards, C., Zhai, C., \& Ji, H. (2022).
Text2Mol: Cross-modal molecule retrieval with natural language queries.
\textit{Proceedings of EMNLP}, 595-607.

\bibitem{radford2021learning}
Radford, A., Kim, J. W., Hallacy, C., et al. (2021).
Learning transferable visual models from natural language supervision.
\textit{International Conference on Machine Learning}, 8748-8763.

\end{thebibliography}

\end{document}

